\documentclass[12pt, a4paper]{article}
\usepackage[utf8]{inputenc}
\usepackage[T2A]{fontenc}
\usepackage[russian]{babel}
\usepackage{amsmath, amsfonts, amssymb, amsthm}
\usepackage[margin=2cm]{geometry}
\usepackage{titlesec}
\usepackage{titletoc}
\usepackage[normalem]{ulem}

% --- НАСТРОЙКА ОГЛАВЛЕНИЯ (titletoc) ---

% Полностью отключаем печать слова "Содержание" и скрытые отступы
\makeatletter
\renewcommand{\tableofcontents}{\@starttoc{toc}}
\makeatother

% Настройка для пунктов программы (I, II, ...)
\titlecontents{section}
  [0em] % Отступ слева
  {\addvspace{0.5em}\bfseries} % Небольшой отступ сверху и жирный шрифт
  {\thecontentslabel\hspace{0.4em}} % Номер, фиксированный пробел 0.4em и сразу текст
  {}
  {\hfill\contentspage} % Номер страницы прижат вправо

% Настройка для подпунктов (1, 2, ...)
\titlecontents{subsection}
  [0.8em] % Сдвинули левее
  {} % Без дополнительных отступов сверху
  {\thecontentslabel\hspace{0.4em}} % Номер, уменьшенный пробел 0.4em и название
  {}
  {\titlerule*[0.5em]{.}\contentspage} % Заполнитель точками до номера страницы


% --- НАСТРОЙКА НУМЕРАЦИИ И ЗАГОЛОВКОВ В ОСНОВНОМ ТЕКСТЕ ---
% Пункты (секции) нумеруются римскими цифрами с точкой
\renewcommand{\thesection}{\Roman{section}.}
% Подпункты (подсекции) нумеруются арабскими цифрами с точкой
\renewcommand{\thesubsection}{\arabic{subsection}.}

% Настройка внешнего вида заголовков
\titleformat{\section}{\normalfont\Large\bfseries}{\thesection}{0.3em}{}
% Уменьшаем отступы у подпунктов
\titlespacing{\subsection}{0pt}{2ex}{1ex}



% ===== Дополнительные пакеты и команды из конспекта =====
\usepackage{tabularx}
\usepackage{epigraph}
\usepackage[dvipsnames]{xcolor}
% Добавлено unicode=true для правильного отображения русского языка в закладках PDF
\usepackage[unicode=true, colorlinks=true, linkcolor=blue, urlcolor=blue]{hyperref}

\DeclareMathOperator{\dif}{\mathsf{d}}
\DeclareMathOperator{\Dif}{\mathsf{D}}
\DeclareMathOperator{\im}{Im}
\DeclareMathOperator{\Ker}{Ker}
\newcommand{\QED}{$\hfill\square$}
\renewcommand{\le}{\leqslant}  % ← было пропущено
\renewcommand{\ge}{\geqslant}  % ← было пропущено
\renewcommand{\phi}{\varphi}   % ← было пропущено
\renewcommand{\L}{\mathcal{L}}   % ← было пропущено
\renewcommand{\subset}{\subseteq}   % ← было пропущено
\newcommand{\weakconv}{\overset{\textit{сл}}{\;\rightharpoonup\;}}

\newcommand{\green}[1]{\textcolor{Green}{#1}}
\newcommand{\red}[1]{\textcolor{BrickRed}{#1}}


\begin{document}

\begin{center}
    \Large \textbf{ЭКЗАМЕНАЦИОННАЯ ПРОГРАММА}\\
    \large по курсу «Функциональный анализ»\\
    3 курс, 6 семестр, 2025/2026 уч.г.\\
    (Поток Коновалова С. П.) \\
    \normalsize
    \vspace{0.8em}
    Авторы: \ \href{https://t.me/andreyko2742}{Комагоров Андрей}, \ Лизюра Дмитрий, \ Candy Club \ | \ \href{https://github.com/LecturesTexClub/lectures_tex_club}{Проект на GitHub}
\end{center}

% Само оглавление
\tableofcontents

\section*{Примечание}

Данный конспект покрывает экзаменационную программу по соответствующему курсу и только её. За основу взят конспект 2025 года и немного из 2022. При этом исправлено всюду плотное множество ошибок, опечаток и неточностей (в том числе текст проверен посредством ИИ). В некоторых местах доказательное изложение отличается от приведённого на лекциях, однако все определения и формулировки сохранены (с точностью до очевидной эквивалентности). Где допустимо, вместо выражений вида $\sup_{\|x\| = 1}\{\}$ супремум берётся по $\|x\| \le 1$, чтобы не оговаривать отдельно случай нульмерного пространства.

\section*{Напоминание}
\addcontentsline{toc}{section}{Напоминание}
 
Данные теоремы из прошлого семестра могут в дальнейшем упоминаться, поэтому для удобства чтения они включены и сюда.

\textbf{Определение.} Непустое подмножество $S$ нормированного пространства $E$ называется \textit{линейным многообразием}, если $S$ замкнуто относительно линейных операций в $E$.

\textbf{Определение.} Непустое подмножество $S$ нормированного пространства $E$ называется \textit{подпространством}, если это замкнутое линейное многообразие в $E$.
 
\phantomsection \label{th:4-1} \textbf{Теорема.} (Рисса, 4.1) Пусть $E$ --- нормированное пространство, $\dim(E) = \infty$.
Тогда единичная сфера в $E$ не является компактом (и даже не вполне ограниченным).
 
\textbf{Следствие.} Единичная сфера в $E$ является компактом тогда и только тогда, когда $\dim(E) < \infty$.
 
\phantomsection \label{th:4-3} \textbf{Теорема.} (4.3, О проекции) Пусть $H$ --- гильбертово пространство, $M \subset H$ --- подпространство. Тогда $H = M \oplus M^\bot$.
 
\phantomsection \label{th:5-3} \textbf{Теорема.} (5.3) Пусть $E_1$ --- нормированное пространство, $E_2$ банахово.
Пусть $S \subset E_1$ --- линейное многообразие, $\overline{S} = E_1$, $A \in \mathcal L(S, E_2)$.
Тогда существует единственный $\widetilde A \in \mathcal L(E_1, E_2)$, такой что $\widetilde A|_S \equiv A$, причём $\|\widetilde A\| = \|A\|$.

\phantomsection \label{th:5-4} \textbf{Теорема.} (5.4, Банаха--Штейнгауза) 
Пусть $E_1$ --- банахово пространство, $E_2$ --- нормированное пространство, $S \subset \mathcal L(E_1, E_2)$, $S \neq \varnothing$.
Если для всех $x \in E_1$ выполнено $\sup_{A \in S} \|Ax\| < \infty$, то $\sup_{A \in S} \|A\| < \infty$.
 
\phantomsection \label{th:5-6} \textbf{Теорема.} (5.6, Критерий поточечной сходимости оператора из $\mathcal L(E_1, E_2)$)
Пусть $E_1$ --- банахово пространство, $E_2$ --- нормированное пространство, $A_n, A \in \mathcal L(E_1, E_2)$.
$A_n \to A$ поточечно тогда и только тогда, когда $\{\|A_n\|\}$ ограничена и $A_n y \to Ay$ для всех $y$ из некоторого подмножества $S \subset E_1$, такого что $\overline{[S]} = E_1$.

\phantomsection \label{th:6-1} \textbf{Теорема.} (6.1, Рисса--Фреше) Пусть $H$ --- гильбертово пространство, $f \in H^*$, тогда существует единственный $y_0 \in H$, такой что $f(x) = (x,y_0)$.
Как следствие, $\|f\|_{H^*} = \|y_0\|_H$.
 
\textbf{The Теорема.} (6.2, Хана--Банаха) Пусть $E$ --- нормированное пространство над $\mathbb K \in \{\mathbb R, \mathbb C\}$, $M \subset E$ --- линейное многообразие, $f$ --- линейный ограниченный функционал, определённый на $M$.
Тогда существует функционал $\widetilde f \in E^*$, такой что $\widetilde f|_M = f$ и $\|\widetilde f\| = \|f\|$.
 
\phantomsection \label{th:hahn-banach-coll-2} \textbf{Следствие 2.} Для всех $x \in E$ найдётся $f \in E^*$, такой что $\|f\| \le 1$ и $f(x) = \|x\|$.
 
\phantomsection \label{th:hahn-banach-coll-3} \textbf{Следствие 3.} Если $f(x) = f(y)$ для всех $f \in E^*$, то $x = y$.
 
\phantomsection \label{th:hahn-banach-coll-4} \textbf{Следствие 4.}
Для всех $x \in E$
\[
    \|x\| = \sup_{\|f\|_{E^*} \le 1} |f(x)|.
\]

\section{(§7) Слабая сходимость в банаховом пространстве}

\subsection{Изометричность вложения \texorpdfstring{$E$}{E} в \texorpdfstring{$E^{**}$}{E**}. Критерий слабой сходимости последовательности.}

\textbf{Определение.} Пусть $E$ --- нормированное пространство.
Последовательность $\{x_n\}$ \textit{слабо сходится} к $x$, если для любого $f \in E^*$ выполнено $f(x_n) \to f(x)$. Обозначение: $x_n \xrightarrow{w} x$.

\textbf{Замечание.} По \hyperref[th:hahn-banach-coll-3]{следствию 3 теоремы Хана--Банаха} слабый предел единственный.

\textbf{Пример.} В гильбертовом пространстве по \hyperref[th:6-1]{теореме Рисса--Фреше} любой функционал представляется в виде $f(x) = (x, y)$.
Из этого следует, что $x_n \xrightarrow{w} x \iff (x_n, y) \to (x, y)~\forall y \in H$.

\textbf{Пример.} В конечномерном пространстве слабая сходимость эквивалентна сходимости по норме.

\textbf{Утверждение.} Если $\|x_n - x\| \to 0$, то $x_n \xrightarrow{w} x$.
Следует из определения нормы: $|f(x_n) - f(x)| \le \|f\| \cdot \|x_n - x\|$.

\textbf{Замечание.} Обратное следствие неверно: рассмотрим $l_2$ и базисные векторы.

\phantomsection \label{th:7-0} \textbf{Утверждение.} (7.0) Пусть $E$ --- нормированное пространство, $x_n \xrightarrow{w} x$. Тогда \\ $\|x\| \le \varliminf_{n \to \infty}\|x_n\|$.

\textbf{Доказательство.} Возьмём волшебный функционал из \hyperref[th:hahn-banach-coll-2]{следствия 2 теоремы Хана--Банаха}: $|f(x)| = \|x\|$, $\|f\| \le 1$. Тогда из слабой сходимости
\[
    \|x\| = |f(x)| = \lim_{n\to\infty}|f(x_n)| \le \varliminf_{n \to \infty}\|f\|\|x_n\| \le \varliminf_{n \to \infty}\|x_n\|.
\]

\QED

\textbf{Упражнение.} Пусть $H$ --- гильбертово пространство, $x_n \xrightarrow{w} x$ и $\|x_n\| \to \|x\|$.
Тогда $\|x_n - x\| \to 0$.

\textbf{Теорема.} (об изометрическом вложении $E$ в $E^{**}$) $\pi: E \to E^{**}, \pi x(f) = f(x)$. Тогда $\forall x \ \|x\| =\|\pi x\|$.

\textbf{Доказательство.}
В силу \hyperref[th:hahn-banach-coll-4]{следствия 4 из теоремы Хана--Банаха}

\[ \| \pi x \| = \sup_{\|f\|\le1} |\pi x(f)|=\sup_{\| f \| \le 1} |f(x)|=\|x\| \]

\QED

\textbf{Замечание.}
    \[x_n \xrightarrow{w} x \Longleftrightarrow \forall f \ f(x_n) \to f(x) \Longleftrightarrow \forall f \ \pi x_n(f) \to \pi x(f) \]
    Это означает поточечную сходимость $F_{x_n}$ к $F_x$.

\phantomsection \label{th:7-1} \textbf{Теорема.} (7.1, критерий слабой сходимости) Пусть $E$ --- нормированное пространство, $\{x_n\} \subset E$.
Тогда $x_n \xrightarrow{w} x$ тогда и только тогда, когда $\{\|x_n\|\}$ ограничена и $f(x_n) \to f(x)$ для любого $f$ из некоторого подмножества $S \subset E^*$, такого что $\overline{[S]} = E^*$.

\textbf{Доказательство.} Условие похоже на \hyperref[th:5-6]{теорему 5.6} (критерий поточечной сходимости), к ней и сведём.
Положим $\pi: E \to E^{**}$, $\pi(x) = F_x$, где $F_x(f) = f(x)$.
По предыдущей теореме $\|x_n\| = \|F_{x_n}\|$, и, как замечено, $x_n \xrightarrow{w} x$ тогда и только тогда, когда $F_{x_n} \to F_x$ поточечно. Теперь возьмём $E^*$ в качестве банахова пространства $E_1$, $\mathbb K$ --- в качестве нормированного пространства $E_2$, что завершает сведение к теореме 5.6.

\QED

\subsection{Слабая сходимость и линейные ограниченные операторы. Слабо ограниченные множества. Теорема Хана.}


\phantomsection \label{th:7-2} \textbf{Теорема.} (7.2) Пусть $E_1, E_2$ --- нормированные пространства, $A \in \mathcal L(E_1, E_2)$.
Если $x_n \xrightarrow{w} x$, то $Ax_n \xrightarrow{w} Ax$.

\textbf{Доказательство.} Рассмотрим $g \in E_2^*$, положим $f = g \circ A$.
Тогда $f(x_n) \to f(x)$ из слабой сходимости $x_n \xrightarrow{w} x$, или, что эквивалентно, $g(Ax_n) \to g(Ax)$, что завершает доказательство.

\QED

\textbf{Определение.}
$S \subset E$~--- слабо ограничено, если $\forall f \in E^* \: f(S)$~--- ограничено

\phantomsection \label{th:7-3} \textbf{Теорема.} (7.3, Хана) Пусть $E$ --- нормированное пространство, $S \subset E$.
$S$ ограничено тогда и только тогда, когда $S$ слабо ограничено.

\textbf{Доказательство.}
$\Rightarrow$ очевидно. $\Leftarrow$ Заметим, что слабая ограниченность $S$ в $E$ эквивалентна поточечной ограниченности $\pi(S)$ в $E^{**}$: для $S = \varnothing$ очевидно, а иначе
\[
    \forall f \in E^* \, \sup_{x \in S} |f(x)| < \infty \ \Longleftrightarrow \ \forall f \in E^* \, \sup_{x \in S} |\pi(x)(f)| < \infty \ \Longleftrightarrow \ \forall f \in E^* \! \sup_{y \in \pi(S)} |y(f)| < \infty.
\]
Тогда по \hyperref[th:5-4]{теореме Банаха-Штейнгауза} $\pi(S)$ ограничено по норме. Следовательно, и $S$ ограничено, так как $\pi$ --- изометрия.

\QED



\subsection{Замкнутый шар в гильбертовом пространстве слабо секвенциально компактен (теорема Банаха).}

\textbf{Лемма.} \phantomsection \label{7.4_lemma} Пусть $H$ --- Гильбертово пространство, $\{y_n\} \subset H$, $\|y_n\| \le R$. Если $(y_n, y)$ сходится для любого $y \in S$, такого что $\overline{[S]} = H$, то $y_n$ слабо сходится к некому $y_0 \in H$.

\textbf{Доказательство.} Для начала введём $S^* = \{(\cdot, y) \mid y \in S\} \subset H^*$ и заметим, что раз $\overline{[S]} = H$, то из \hyperref[th:6-1]{теоремы Рисса--Фреше} следует, что $\overline{[S^*]} = H^*$. Но прежде чем применить \hyperref[th:7-1]{критерий слабой сходимости}, нужно понять, к чему сходится $(y_n, y)$. Для этого рассмотрим линейный функционал $f(y) = \lim_{n \to \infty} (y, y_n)$ на $[S]$. Он ограничен, так как $|f(y)| \le \|y\|\sup_{n} \|y_n\| \le R\|y\|$. Значит, его можно продолжить до линейного функционала $F$ на $\overline{[S]} = H$ с сохранением нормы. Тогда по теореме Рисса--Фреше найдётся $y_0 \in H$, такой что $F(y) = (y, y_0)$ для всех $y \in H$. На $S$ мы знаем, что $F(y) = f(y) = \lim_{n \to \infty} (y, y_n) = (y, y_0)$. Вот теперь по критерию слабой сходимости получаем требуемое. 

\QED

\phantomsection \label{th:7-4} \textbf{Теорема.} (7.4, Банаха) Замкнутый шар $\overline B(0, R)$ в гильбертовом пространстве слабо секвенциально компактен.

\textbf{Доказательство.} Возьмём последовательность $\{x_n\} \subset \overline B(0, R)$ и попробуем найти подпоследовательность $\{y_n\}$, слабо сходящуюся к $y \in \overline B(0, R)$ (это и означает слабую секвенциальную компактность). Причём $\|y\| \le R$ выполнено автоматически по \hyperref[th:7-0]{утверждению 7.0}.

Поступим хитро: рассмотрим $L = \overline{[\{x_n\}]}$ и выделим $\{y_n\}$, слабо сходящуюся в $L$. По \hyperref[th:4-3]{теореме о проекции} $H = L \oplus L^\bot$, и тогда для любого $v \in H$ получим $(y_n, v) = (y_n, l) + (y_n, l^\bot) = (y_n, l) + 0$ из ортогональности. То есть получится слабая сходимость во всём $H$. 

Чтобы применить лемму, нам нужна ограниченность $\{\|y_n\|\}$ --- выполнено автоматически --- и сходимость $(y_n, y)$ на множестве $S$, таком что $\overline{[S]} = L$. Таковым будет $\{x_n\}$.

Выделим $\{y_n\}$ диагональным методом Кантора.
Рассмотрим последовательность \\ $\{(x_n, x_1)\}_{n=1}^\infty$.
Она ограничена, значит, есть сходящаяся $\{(x_n^1, x_1)\}$ по теореме Больцано--Вейерштрасса.
Теперь рассмотрим последовательность $\{(x_n^1, x_2)\}$, аналогично выделяется сходящаяся $\{(x_n^2, x_2)\}$.
Продолжаем так делать для всех элементов, после чего возьмём последовательность $\{y_n := x_n^n\}$.
Так как для всех $m$ последовательность ${(y_n, x_m)}_{n = 1}^\infty$ сходится, то по лемме $y_n$ слабо сходится.



\QED


\section{(§8) Обратный оператор. Обратимость}

\subsection{Обратимость линейного, ограниченного снизу (\texorpdfstring{$\|Ax\| \ge k\|x\|$}{||Ax|| >= k||x||}) оператора.}

Пусть $E_1, E_2$ --- нормированные пространства.

\textbf{Определение.} (загадочное) Оператор $A: E_1 \to E_2$ называется \textit{обратимым}, если он инъективен. Оператор $A^{-1}: \im A \to E_1$ называется \textit{обратным} к $A$, если $A^{-1}A = I_{E_1}$ и $AA^{-1} = I_{\im A}$.

\textbf{Определение.} $A$ называется \textit{ограниченным снизу}, если найдётся $m > 0$, такое что $\|A x\| \ge m \|x\|$ для всех $x$ (или, эквивалентно, $\|x\| \le M\|Ax\|$ для некоторого $M$).

\phantomsection \label{th:8-1} \textbf{Теорема.} (8.1) Пусть $E_1, E_2$ --- нормированные пространства, $A \in \mathcal L(E_1, E_2)$.
Оператор $A^{-1}$ существует и ограничен на $\im A$ тогда и только тогда, когда $A$ ограничен снизу.

\textbf{Доказательство.} $\Rightarrow$. Пусть $x \in E_1$, тогда $\|x\| = \|A^{-1} A x\| \le \|A^{-1}\| \cdot \|Ax\| =: M\|Ax\|$.

$\Leftarrow$. Заметим, что ядро тривиально, следовательно, $A$ обратим.
Теперь докажем ограниченность обратного: пусть $y = Ax$, тогда
\[
    \|A^{-1} y\| = \|x\| \le M\|Ax\| = M\|y\| \Longrightarrow \|A^{-1}\| \le M.
\]

\QED



\subsection{Обратимость возмущённого оператора \texorpdfstring{$A + \Delta A$}{A + Delta A}.}

\phantomsection \label{th:8-2} \textbf{Теорема.} (8.2) Пусть $E$ --- банахово пространство, $A \in \mathcal L(E)$, $\|A\| < 1$.
Тогда $I + A$ обратим и $(I + A)^{-1} \in \mathcal L(E)$.

\textbf{Доказательство.} Идея: если смотреть на $I$ и $A$ как на числа, то получится
\[
    \frac{1}{I + A} = \frac{1}{1 + A} = \sum_{k=0}^{\infty} (-A)^k.
\]
Так и сделаем: рассмотрим ряд $S = \sum_{k=0}^{\infty} (-A)^k$, который сходится в $\mathcal{L}(E)$ за счёт $\|A\| < 1$. Пусть $S_n$ --- его конечная сумма. 
Тогда 
\[
    S_n(I + A) = (I + A) S_n = I + (-1)^n A^{n+1} \to I.
\]
Следовательно,
\[
    S(I + A) = (I + A)S = I.
\]
Значит, $S = (I + A)^{-1}$.

\QED

\phantomsection \label{th:8-3} \textbf{Теорема.} (8.3) Пусть $E$ --- банахово пространство, $A$ обратим, $A, A^{-1}, \Delta A \in \mathcal L(E)$.
Если $\|A^{-1}\|\|\Delta A\| < 1$, то $(A + \Delta A)^{-1} \in \mathcal L(E)$.
Очевидно из теоремы 8.2, т.к. композиция обратимых обратима:
\[
    A + \Delta A = A (I + A^{-1} \Delta A).
\]



\subsection{Формулировка теоремы Банаха об обратном операторе. Доказательство теоремы Банаха в случае гильбертова пространства.}

\textbf{Лемма.} (8.4) \phantomsection \label{8.4_lemma}
        Пусть $E$ --- банахово пространство,
        $A \in \L(E)$ --- ограниченный снизу оператор. 
        Тогда $\overline{\im A} = \im A$.

\textbf{Доказательство.}
    Пусть \(\left\{y_n \right\} \subset \im A\), такая что \(y_n \rightarrow y\) по норме. Представим как $y_n = A z_n, z_n \in E$. Последовательность $\{y_n\}$ сходится, а значит она фундаментальна. Из ограниченности снизу:
    \[ \varepsilon > \|y_{n+p}-y_n\|=\|Az_{n+p}-Az_n\| \ge m \cdot \|z_{n+p} - z_n \|.\] 
    Следовательно, $\{z_n\}$ - фундаментальна в $E$~--- банаховом, откуда $\exists \lim z_n = z$. Получаем $y_n = A z_n \to A z = y$. То есть $y \in \im A$, а значит $\im A$ замкнут.

\QED

\phantomsection \label{th:8-4} \textbf{Теорема.} (8.4, Банаха об обратном операторе) Пусть $E_1, E_2$ --- банаховы пространства, $A \in \mathcal L(E_1, E_2)$ --- биекция.
Тогда $A^{-1} \in \mathcal L(E_2, E_1)$.

\textbf{Доказательство.} Будем доказывать для $E_1 = E_2 = H$ гильбертовых.

\begin{enumerate}
    \item Покажем, что $A^* \colon H \to H$ обратим (где $A^*$ --- эрмитово сопряжённый оператор, см. §10).
    По \hyperref[th:10-2]{теореме 10.2} $\overline{\im A} \oplus \Ker A^* = H$. По условию $\im A = H$, откуда $\Ker A^* = \{0\}$, то есть $A^*$ обратим.

    \item Докажем, что $(A^*)^{-1} \in \mathcal{L}(H)$. По \hyperref[th:8-1]{теореме 8.1} для этого достаточно показать, что
    $A^*$~--- ограниченный снизу, то есть 
    $\exists M > 0\colon \forall x \ \|x\| \le M\|A^* x\|$. С учётом обратимости~$A^*$ это эквивалентно условию $\exists M > 0\colon \forall x \ (\|A^* x\| = 1 \rightarrow \|x\| \le M)$.
    То есть надо доказать, что множество $S := \{x \in H \mid \|A^*x\| = 1\}$ ограничено.
    
    Покажем, что $S$~--- слабо ограниченное, то есть
    \[
    \forall y \in H\; \exists M_y\; \forall x \in S \left(|(x, y)| \le M_y \right)
    \]
    тогда по \hyperref[th:7-3]{теореме Хана} будет следовать ограниченность $S$. (Здесь мы пользуемся \hyperref[th:6-1]{теоремой Рисса--Фреше}, заменяя произвольный линейный функционал на скалярное произведение с элементом $y \in H$.)

    Используем, что $A$~--- биекция $H$ на $H$: $\forall y \in H \; \exists z \in H \colon Az = y$

    \[ |(x, y)| = |(x, Az)| = |(A^*x, z)| \le \underbrace{\|A^* x\|}_{=1\text{ на $S$}} \cdot \underbrace{\|z\|}_{:=M_y}\]
    Действительно слабо ограничено, а значит и просто ограничено и оператор $A^*$ ограничен снизу.

    \item Докажем, что $A^*$ --- биекция
    По \hyperref[th:10-2]{теореме 10.2} имеем 
    $\overline{\im A^*} \oplus \Ker A = H$. Следовательно, $\overline{\im A^*} = H$ (так как $\ker A = \{0\}$ по условию). Значит, по \hyperref[8.4_lemma]{лемме 8.4} $\im A^* = \overline{\im A^*} = H$. Имеем инъективность (по пункту 1) и сюрьективность, то есть $A^*$~--- биекция.

    \item Выходит что $A^*$ удовлетворяет условиям теоремы и $(A^*)^{-1} \in \L(H)$.
    Провернём рассуждения ещё раз для $A^*$, получим $(\underbrace{A^{**}}_{=A})^{-1} \in \L(H)$~--- перешли к исходному оператору $A$. 
\end{enumerate}

\QED



\section{(§9) Спектр. Резольвента}

\subsection{Операторнозначные функции комплексного переменного. Аналитичность резольвенты. Спектральный радиус. Основная теорема о спектре.}

\textbf{Напоминание.} Аналог данной темы был на алгеме: для оператора $A$ мы рассматриваем решения уравнения $(A - \lambda I) x = 0$ и $(A - \lambda I) x = y$.
Иными словами, исследуем ядро оператора $(A - \lambda I)$ в первом случае, а во втором --- его образ.

В данном параграфе мы будем активно применять методы комплексного анализа, поэтому везде будем считать, что $E$ --- банахово пространство над полем $\mathbb C$.
Также для оператора $A \in \mathcal L(E)$ и $\lambda \in \mathbb C$ мы будем обозначать $A_\lambda = A - \lambda I$.

\textbf{Определение.} Пусть $A \in \mathcal L(E)$.
\begin{itemize}
    \item \textit{Резольвентным множеством} называется $\rho(A) = \{\lambda \in \mathbb C~|~A_\lambda^{-1} \in \mathcal L(E)\}$, то есть ``хорошее`` множество, на котором всё обратимо.
    \item \textit{Спектром} называется $\sigma(A) = \mathbb C \setminus \rho(A)$, то есть ``плохое`` множество.
    \item \textit{Резольвентой} для $\lambda \in \rho(A)$ называется оператор $R_\lambda := A_\lambda^{-1}$.
    \item \textit{Точечным спектром}, или \textit{собственными значениями}, называется $\sigma_p(A)$ --- множество таких $\lambda \in \mathbb C$, что уравнение $(A - \lambda I) x = 0$ имеет нетривиальное решение, называемое \textit{собственным вектором}.
    \item \textit{Непрерывным спектром} называется $\sigma_c(A)$ --- множество таких $\lambda \in \mathbb C$, что $\Ker A_\lambda = \{0\}$, но $\im A_\lambda \ne E$ и $\overline{\im A_\lambda} = E$.
    \item \textit{Остаточным спектром} называется $\sigma_r(A)$ --- множество таких $\lambda \in \mathbb C$, что $\Ker A_\lambda = \{0\}$, но $\overline{\im A_\lambda} \ne E$.
\end{itemize}

\textbf{Замечание.} По \hyperref[th:8-4]{теореме Банаха об обратном операторе} $\sigma(A) = \sigma_p(A) \sqcup \sigma_c(A) \sqcup \sigma_r(A)$. При $\dim E < \infty$ выполнено $\sigma(A) = \sigma_p(A)$, а иначе --- необязательно.

\phantomsection \label{th:9-1} \textbf{Теорема.} (9.1) Пусть $A \in \mathcal L(E)$, $E \neq \{0\}$.
Тогда
\begin{enumerate}
    \item $\sigma(A)$ --- непустое замкнутое множество.
    \item Пусть $r(A) := \sup_{\lambda \in \sigma(A)} |\lambda|$, тогда $r(A) = \lim_{n \to \infty} \sqrt[n]{\|A^n\|}$.
\end{enumerate}

Ниже приведено доказательство, использующее ТФКП для функций из $\mathbb{C}$ в $\mathcal{L}(E)$. Утверждается, что все нужные нам теоремы (сходимость рядов Лорана и теорема Лиувилля) естественным образом переносятся на этот случай. Однако можно проделать аналогичные рассуждения, используя только теорию для функций из $\mathbb{C}$ в $\mathbb{C}$. Например, \href{https://proofwiki.org/wiki/Gelfand%27s_Spectral_Radius_Formula/Banach_Algebra}{здесь} изложено доказательство формулы для $r(A)$ (только здесь $A$ --- это $\mathcal{L}(E)$, unital Banach algebra, а $x$ --- это оператор). По сути всё сводится к рассмотрению композиции $R: \lambda \mapsto R_\lambda$ с произвольным линейным непрерывным функционалом $\phi: \mathcal{L}(E) \to \mathbb{C}$. При этом нужно в любом случае использовать, что $R$ дифференцируемо и что $\lambda^n \in \sigma(A^n)$ (то есть все пункты ниже, кроме 4-5). Для непустоты спектра приём аналогичный: в предположении противного имеем, что для каждого $\phi$ функция $\phi R$ целая и по теореме Лиувилля равна нулю. Значит, $R = 0$ по \hyperref[th:hahn-banach-coll-4]{следствию 4 из теоремы Хана--Банаха}.

\textbf{Доказательство.}  Заметим, что $R_\lambda$ можно рассматривать, как оператор $R: \rho(A) \to \mathcal L(E)$.

Далее доказательство разбивается на 9 пунктов:
\begin{enumerate}
    \setcounter{enumi}{-2}
    \item Если $|\lambda| > \|A\|$, то $\lambda \in \rho(A)$.
        Следствие из \hyperref[th:8-2]{теоремы 8.2}.

    \item $\rho(A)$ --- открытое множество.
        Следствие из \hyperref[th:8-3]{теоремы 8.3}, что доказывает замкнутость спектра (следовательно, компактность).

    \item $\lambda \mapsto R_\lambda$ непрерывно на $\rho(A)$.
        Аналогично доказательству теоремы 8.2 рассмотрим $B = A_{\lambda_0}$, $\Delta B = -\lambda I$ для достаточно малых $\lambda$, тогда
        \[
            (B + \Delta B)^{-1} = (B(I + B^{-1} \Delta B))^{-1} = (I + B^{-1} \Delta B)^{-1}B^{-1} = \sum_{k=0}^{\infty} (-1)^k  (B^{-1} \Delta B)^k B^{-1}.
        \]
        Теперь
        \[
            (B + \Delta B)^{-1} - B^{-1} = \sum_{k=1}^{\infty} (-1)^k (B^{-1} \Delta B)^k B^{-1},
        \]
        значит,
        \[
            \|(B + \Delta B)^{-1} - B^{-1}\| \le \sum_{k=1}^{\infty} \|B^{-1}\|^{1 + k} \cdot \|\Delta B\|^k = \frac{\|B^{-1}\|^2 \cdot \|\Delta B\|}{1 - \|B^{-1}\| \cdot \|\Delta B\|}.
        \]
        Обратно к $A_\lambda$:
        \[
            \|R_{\lambda_0 + \lambda} - R_{\lambda_0}\| \le \frac{\|R_{\lambda_0}\|^2 \cdot |\lambda|}{1 - \|R_{\lambda_0}\| \cdot |\lambda|}.
        \]
        При $\lambda \to 0$ это стремится к нулю, что доказывает непрерывность.

    \item Равенство Гильберта: $R_\lambda - R_{\lambda_0} = (\lambda - \lambda_0) R_\lambda R_{\lambda_0}$.
        Заметим, что
        \[
            A_{\lambda} (R_\lambda - R_{\lambda_0}) A_{\lambda_0} = A_{\lambda_0} - A_{\lambda} = (\lambda - \lambda_0) I.
        \]
        Остаётся домножить на $R_\lambda$ слева и на $R_{\lambda_0}$ --- справа.

    \item $\lambda \mapsto R_\lambda$ дифференцируема на $\rho(A)$ (см. главу про нелинейный анализ). Действительно, по равенству Гильберта:
        \[
            R_\lambda - R_{\lambda_0} = (\lambda - \lambda_0) R_{\lambda_0}^2 + (\lambda - \lambda_0) (R_\lambda - R_{\lambda_0}) R_{\lambda_0},
        \]
        где $(\lambda - \lambda_0) (R_\lambda - R_{\lambda_0}) R_{\lambda_0} = o(|\lambda - \lambda_0|)$ по непрерывности $R_\lambda$.

    \item $\sigma(A) \ne \varnothing$.
        От противного: $\rho(A) = \mathbb C$, тогда $R_\lambda$ целая.
        Отсюда при $|\lambda| > \|A\|$ выполнено
        \[
            \|R_\lambda\| \le \frac{1}{|\lambda|} \cdot \frac{1}{1 - \frac{\|A\|}{|\lambda|}},
        \]
        что стремится к нулю при $\lambda \to \infty$, значит, $R_\lambda = 0$ по теореме Лиувилля из ТФКП.
        Противоречие: нет обратного оператора.
        
    \item $r(A)$ --- это радиус сходимости \textit{ряда Неймана} (Не Фона Неймана, а Карла Неймана):
        \[
            R_{\lambda} = -\frac{1}{\lambda} \sum_{n=0}^\infty \frac{1}{\lambda^n} A^n.
        \]
        Мы знаем, что равенство верно при $|\lambda| > \|A\|$, а также, что $R$ дифференцируема на $\rho(A)$.
        Рассмотрим случаи:
        \begin{itemize}
            \item $|\lambda_0| > r(A)$, докажем сходимость.
                Действительно, здесь $R$ голоморфна, а значит, раскладывается в ряд Лорана, причём единственным образом, и разложение мы уже знаем.

            \item $0 < |\lambda_0| < r(A)$, докажем расходимость.
                От противного, тогда ряд бы сходился и при всех бóльших по модулю $\lambda$, ибо это ряд Лорана в бесконечности.
                Но из-за этого получается, что $r(A) \le |\lambda_0|$, ибо данный ряд был бы обратным к $A_\lambda$:
                \[
                    (A - \lambda I) S_n = S_n (A - \lambda I) = I - \frac{A^{n+1}}{\lambda^{n+1}},
                \]
                где $S_n = -\frac{1}{\lambda} \sum_{k=0}^n \frac{1}{\lambda^k} A^k$, причём $\frac{A^{n+1}}{\lambda^{n+1}} \to 0$ как член сходящегося ряда.
        \end{itemize}

    \item Если $\lambda \in \sigma(A)$, то $\lambda^n \in \sigma(A^n)$.
        (Аналог факта из линейной алгебры про собственные значения матрицы)

        От противного: допустим, что $\lambda^n \in \rho(A^n)$, то есть $(A^n - \lambda^n I)^{-1} \in \mathcal L(E)$.
        Распишем:
        \[
            A^n - \lambda^n I = (A - \lambda I)(A^{n-1} + \lambda A^{n-2} + \dots + \lambda^{n-1} I).
        \]
        Применим справа к обеим частям оператор $(A^n - \lambda^n I)^{-1}$, тогда слева будет $I$, а справа --- композиция $(A - \lambda I)$ и какого-то оператора, то есть правого обратного.
        Но, так как степенные операторы перестановочны, можно аналогично получить, что есть и левый обратный, откуда $(A - \lambda I)$ обратим --- противоречие.

    \item $r(A) = \lim_{n \to \infty} \sqrt[n]{\|A^n\|}$.
        В пункте 5 мы уже доказали, что $r(A) = \overline{\lim}_{n \to \infty} \sqrt[n]{\|A^n\|}$, осталось убрать верхний предел.
        Заметим, что по пункту 6
        \[
            r(A) \le \sqrt[n]{r(A^n)} \le \sqrt[n]{\|A^n\|} =: \alpha_n.
        \]
        Значит, $r(A) \le \alpha_n$ для всех $n$, то есть $\underline{\lim}_{n \to \infty} \alpha_n \ge r(A) = \overline{\lim}_{n \to \infty} \alpha_n$, то есть верхний и нижний предел совпадают, и они равны $r(A)$.
\end{enumerate}

\QED

\section{(§10) Сопряженный оператор}

\subsection{Норма сопряженного оператора (в линейном нормированном пространстве).}

\textbf{Определение.} Пусть $E_1, E_2$ --- два нормированных пространства над полем $\mathbb K$, $A \in \mathcal L(E_1, E_2)$.
Если $g \in E_2^*$, то $x \mapsto g(Ax)$ --- это линейный функционал, действующий из $E_1$ в $\mathbb K$.
В связи с этим можно определить $A^*: E_2^* \to E_1^*$, $(A^* g) x = g(Ax)$.

\phantomsection \label{th:10-1} \textbf{Теорема.} (10.1) $A^* \in \mathcal L(E_2^*, E_1^*)$, а также $\|A^*\| = \|A\|$.

\textbf{Доказательство.} Линейность очевидна.
Ограниченность:
\[
    |(A^* g)(x)| = |g(Ax)| \le \|g\| \cdot \|A\| \cdot \|x\|,
\]
значит,
\[
    \|A^* g\| \le \|A\| \cdot \|g\|,
\]
то есть $\|A^*\| \le \|A\|$.
Теперь докажем обратное неравенство: по \hyperref[th:hahn-banach-coll-4]{следствию 4 теоремы Хана--Банаха}
\[
    \|Ax\| = \sup_{\|g\| \le 1} |g(Ax)| = \sup_{\|g\| \le 1} |(A^* g) x| \le \sup_{\|g\| \le 1} \|A^*\| \cdot \|g\| \cdot \|x\| = \|A^*\| \cdot \|x\|,
\]
то есть $\|A\| \le \|A^*\|$.

\QED



\subsection{Сопряженные операторы в гильбертовом пространстве. Равенство \texorpdfstring{$H = \overline{Im A} \oplus Ker A^*$}{H = Im A (+) Ker A*}.}

Далее мы будем рассматривать только гильбертовы пространства, ибо здесь всё гораздо проще. Мы принимаем соглашение, что скалярное произведение линейно по первому аргументу и сопряжённо-линейно по второму.

\textbf{Определение.} Пусть $H_1, H_2$ --- гильбертовы пространства, $A \in \mathcal L(H_1, H_2)$.
\textit{Эрмитово сопряжённым оператором} $A^*: H_2 \to H_1$ называется такой оператор, что для всех $x \in H_1$, $y \in H_2$ выполнено $(Ax, y) = (x, A^* y)$.

\textbf{Утверждение.} Для любого $A \in \mathcal{L}(H_1, H_2)$ существует единственный эрмитово сопряжённый $A^*$ и $A^* \in \mathcal{L}(H_2, H_1)$.

\textbf{Доказательство.} Единственность следует из \hyperref[th:6-1]{теоремы Рисса--Фреше} и \hyperref[th:hahn-banach-coll-3]{следствия 3 теоремы Хана--Банаха}. Докажем существование. Возьмём $\pi_1 \colon H_1 \to H_1^*$ и $\pi_2 \colon H_2 \to H_2^*$ из теоремы Рисса--Фреше, то есть
\[
    \forall x \in H_1 \ \pi_1 x = (\cdot, x), \quad \forall y \in H_2 \ \pi_2 y = (\cdot, y).
\]
Обозначим эрмитово сопряжённый оператор как $A^*$, а просто сопряжённый как $A'$. Покажем, что $A^* = \pi_1^{-1}A'\pi_2$, то есть
\[
    \forall x \in H_1, y \in H_2 \ (Ax, y) = (x, \pi_1^{-1}A'\pi_2 y).
\]
Это делается по определению $\pi_1$, $A'$ и $\pi_2$:
\[
    (x, \pi_1^{-1}A'\pi_2 y) = (\pi_1(\pi_1^{-1}A'\pi_2 y))x = (A'(\pi_2y))x = (\pi_2y)(Ax) = (Ax, y).
\]

\QED

Далее, говоря о Гильбертовых пространствах, под сопряжённым оператором будем подразумевать именно эрмитово сопряжённый.

\textbf{Упражнение.}
\begin{itemize}
    \item $(\alpha A + \beta B)^* = \overline \alpha A^* + \overline \beta B^*$.
    \item $(AB)^* = B^* A^*$.
    \item $(A^*)^* = A$.
    \item $I^* = I$.
    \item Если $A$ --- биекция, то $(A^*)^{-1} = (A^{-1})^*$.
\end{itemize}

\phantomsection \label{th:10-2} \textbf{Теорема.} (10.2) Пусть $H$ --- гильбертово пространство, $A \in \mathcal L(H)$.
Тогда $\overline{\im A} \oplus \Ker A^* = H$.

\textbf{Доказательство.} Докажем, что $(\im A)^\bot = \Ker A^*$.
Действительно, $y \in (\im A)^\bot \Leftrightarrow$ для всех $x \in H$ выполнено $(Ax, y) = 0 \Leftrightarrow$ для всех $x \in H$ $(x, A^* y) = 0 \Leftrightarrow A^* y = 0 \Leftrightarrow y \in \Ker A^*$.
Теперь по \hyperref[th:4-3]{теореме 4.3} выполнено $\overline{\im A} \oplus \overline{\im A}^\bot = H$, остаётся заметить, что $(\im A)^\bot = (\overline{\im A})^\bot$.

\QED



\section{(§11) Самосопряженные операторы}

\subsection{Свойства квадратичной формы \texorpdfstring{$(Ax, x)$}{(Ax, x)} и собственных значений самосопряженного оператора \texorpdfstring{$A$}{A}.}

В данном параграфе мы будем работать с \textbf{гильбертовым} пространством $H$ над полем $\mathbb C$ по умолчанию.

\textbf{Определение.} Пусть $A \in \mathcal L(H)$.
Он называется \textit{самосопряжённым} (ССО), если $A = A^*$.

\textbf{Утверждение 1.} Если $A$ --- ССО, $M$ --- инвариантное относительно $A$ подпространство, то $M^\bot$ тоже инвариантно относительно $A$.

\textbf{Доказательство.} Если $y \in M^\bot$, то для всех $x \in M$ выполнено $(Ax, y) = 0$ или же $(x, Ay) = 0$, то есть $Ay \in M^\bot$.

\QED

Как проверить, является ли оператор $A$ самосопряжённым?
Один способ --- подставить в скалярное произведение, но это не всегда легко делается.
Ещё один --- рассмотреть квадратичную форму $K_A: H \to \mathbb C$, $K_A(x) = (Ax, x)$.

\textbf{Утверждение 2.} Если $H$ построено над полем $\mathbb C$, то из $K_A = K_B$ следует $A = B$.
(Но над $\mathbb R$ это уже неверно)

\textbf{Доказательство.} Рассмотрим $C = A - B$, тогда для всех $x$ выполнено $(Cx, x) = 0$.
Тогда для всех $x, y \in H$ выполнено
\[
    0 = (C(x + y), x + y) = (Cx, x) + (Cx, y)+ (Cy, x) + (Cy, y) = (Cx, y)+ (Cy, x).
\]
Подставим $x$, $iy$:
\[
    0 = (C(x + iy), x + iy) = (Cx, iy) + (Ciy, x) = -i(Cx, y) + i(Cy, x).
\]
Итого получили, что $(Cx, y) + (Cy, x) = (Cx, y) - (Cy, x) = 0$, то есть $(Cx, y) = 0$. Следовательно, для всех $x$ $Cx = 0$, следовательно, $C = 0$.

\QED

\phantomsection \label{th:11-1} \textbf{Теорема.} (11.1, о свойствах ССО) Пусть $A$ --- ССО в $H$.
Тогда
\begin{enumerate}
    \item Для всех $x \in H$ выполнено $K_A(x) \in \mathbb R$.
        Причём это критерий самосопряжённости.
    \item Если $\lambda$ --- собственное значение $A$, то $\lambda \in \mathbb R$.
    \item Если $e_1, e_2$ --- два собственных вектора, соответствующих разным собственным значениям, то $(e_1, e_2) = 0$.
\end{enumerate}

\textbf{Доказательство.}
\begin{enumerate}
    \item Из свойств скалярного произведения и самосопряжённости получаем
        \[
            K_A(x) = (Ax, x) = \overline{(x, Ax)} = \overline{(Ax, x)} = \overline{K_A(x)}.
        \]
        Обратно: по утверждению 2 достаточно проверить, что $K_{A^*} = K_A$.
        По определению $(A^* x, x) = (x, Ax)$, а $(Ax, x) = \overline{(x, Ax)}$.
        Но $(Ax, x) \in \mathbb R$, поэтому они равны.

    \item Пусть $Ae = \lambda e$, $e \ne 0$. Тогда $\lambda (e, e) = (Ae, e) = (e, Ae) = \overline \lambda (e, e)$, следовательно, $\lambda = \overline \lambda$.

    \item $\lambda_1(e_1, e_2) = (Ae_1, e_2) = (e_1, Ae_2) = \overline \lambda_2 (e_1, e_2)$, откуда по пункту 2 $(e_1, e_2) = 0$.
\end{enumerate}

\QED



\subsection{Разложение гильбертова пространства \texorpdfstring{$H = \overline{Im(A - \lambda I)} \oplus Ker(A - \lambda I)$}{H = Im(A - lambda I) (+) Ker(A - lambda I)}, где \texorpdfstring{$A$}{A} --- самосопряженный оператор.}

\phantomsection \label{th:11-2} \textbf{Теорема.} (11.2) Если $A$ --- ССО, то справедлива формула $\overline{\im A_\lambda} \oplus \Ker A_\lambda = H$.

\textbf{Доказательство.} Это не так просто следует из \hyperref[th:10-2]{теоремы 10.2}, она нам даёт только $\overline{\im A_\lambda} \oplus \Ker (A_\lambda)^* = H$, а $(A_\lambda)^* = A - \overline \lambda I = A_{\overline \lambda}$.
Засим рассмотрим два случая:
\begin{itemize}
    \item $\lambda \in \mathbb R$. Тогда $A_\lambda = A_{\overline \lambda}$.
    \item $\lambda \notin \mathbb R$.
        Из теоремы \hyperref[th:11-1]{теоремы 11.1} получаем, что $\lambda$ и $\overline \lambda$ --- не собственные значения, тогда $\Ker A_\lambda = \Ker A_{\overline \lambda} = \{0\}$.
\end{itemize}

\QED

\textbf{Следствие.} Если $A$ --- ССО, то $\sigma_r(A) = \varnothing$.



\subsection{Критерий принадлежности числа спектру самосопряженного оператора. Вещественность спектра самосопряженного оператора.}

\textbf{Напоминание.} Оператор $A$ \textit{ограничен снизу}, если найдётся $m > 0$, такое что $\|A x\| \ge m \|x\|$ для всех $x$.

\phantomsection \label{th:11-3} \textbf{Теорема.} (11.3) Пусть $A$ --- ССО.
Тогда:
\begin{enumerate}
    \item $\lambda \in \rho(A)$ тогда и только тогда, когда $A_\lambda$ ограничен снизу.
    \item $\lambda \in \sigma(A)$ тогда и только тогда, когда существует $\{x_n\}$, такое что все $\|x_n\| = 1$ и $\|A_\lambda x_n\| \to 0$, то есть ``почти собственное значение``.
\end{enumerate}

\textbf{Доказательство.} Достаточно доказать первый пункт, так как второй получается по контрапозиции. Вспоминая \hyperref[th:8-1]{теорему 8.1}, получаем, что слева направо очевидно.
Справа налево: есть ограниченность снизу, нужно $\lambda \in \rho(A)$.
У нас есть $\overline{\im A_\lambda} \oplus \Ker A_\lambda = H$, причём по \hyperref[8.4_lemma]{лемме 8.4} образ замкнут. $\Ker A_\lambda = \{0\}$ за счёт ограниченности снизу, значит $\im A_\lambda = H$, значит $\lambda \in \rho(A)$.

\QED

\phantomsection \label{th:11-4} \textbf{Теорема.} (11.4) Если $A$ --- ССО в $H$, то $\sigma(A) \subset \mathbb R$.

\textbf{Доказательство.} Пусть $\lambda = \mu + i \nu \in \sigma(A)$.
Докажем, что $\|A_\lambda x\| \ge |\nu| \cdot \|x\|$:
\[
    \|A_\lambda x\|^2 = \|A_\mu x - i \nu x\|^2 = 
\]
\[
    = \|A_\mu x\|^2 + (A_\mu x, - i\nu x) + (-i \nu x, A_\mu x) + |\nu|^2 \|x\|^2.
\]
Итак, $\|A_\mu x\|^2 \ge 0$, а второе и третье слагаемые сокращаются за счёт самосопряжённости $A_\mu$.
Следовательно, по \hyperref[th:11-3]{теореме 11.3} получаем, что $\nu = 0$.

\QED



\subsection{Теорема о спектре самосопряжённого оператора: \texorpdfstring{$\sigma(A) \subseteq [m_-, m_+], r(A) = \|A\|$}{sigma(A) <= [m-, m+], r(A) = ||A||}.}

\textbf{Лемма.} Пусть $A$ --- ССО. Тогда $\|A\| = \sup_{\|x\| \le 1} |(Ax, x)|$.

\textbf{Доказательство.} Запишем определение нормы по следствию 4 теоремы Хана--Банаха: 
\[
    \|A\| = \sup_{\|x\| \le 1} \|Ax\| = \sup_{\|x\| \le 1, \|y\| \le 1} |(Ax, y)|.
\]
Обозначим $R := \sup_{\|x\| \le 1} |(Ax, x)|$. Нужно взять произвольные $x, y$ из единичного шара и доказать, что $|(Ax, y)| \le R$. При этом над полем $\mathbb{C}$ может оказаться, что $\mathrm{Im}(Ax, y) \neq 0$. Тогда домножим на такое $\alpha \in \mathbb{C}$, что $|\alpha| = 1$ и $\mathrm{Im} (\alpha(Ax, y)) = 0$. Получим $|(Ax, y)| = |\alpha(Ax, y)| = |(Ax, \overline\alpha y)|$ и $\|\overline\alpha y\| \le 1$. То есть можно считать, что $(Ax, y)$ вещественно.

Теперь можно перейти от $(Ax, y)$ к $(Az, z)$ следующим образом:
\[
    (Ax, y) \overset{\text{ССО}}{=} \frac{(Ax, y) + (x, Ay)}{2} = \frac{(Ax, y) + (Ay, x)}{2} = \frac{1}{4}\big((A(x+y), x+y) - (A(x-y), x-y)\big).
\]
А дальше нам поможет равенство параллелограмма:
\[
    |(Ax, y)| = \frac{1}{4}\big|(A(x+y), x+y) - (A(x-y), x-y)\big| \le \frac{1}{4}\big(R\|x+y\|^2 + R\|x-y\|^2\big) = \frac{R}{2}(\|x\|^2 + \|y\|^2) \le R.
\]
Этим доказано, что $\|A\| \le R$. Обратное неравенство очевидно.

\QED

\phantomsection \label{th:11-5} \textbf{Теорема.} (11.5) Пусть $H \neq \{0\}$ --- комплексное, $A$ --- ССО.
Тогда:
\begin{enumerate}
    \item $\sigma(A) \subset [m, M]$, где $m = \inf (Ax, x)$, $M = \sup (Ax, x)$, где супремум и инфимум берутся по $\|x\| = 1$.
    \item $m, M \in \sigma(A)$.
    \item $\|A\| = \max(|m|, |M|)$.
\end{enumerate}
Следовательно, спектральный радиус равен $\|A\|$.

\textbf{Доказательство.}
\begin{enumerate}
    \setcounter{enumi}{2}
    \item Следует из леммы:
    \[
        \|A\| = \sup_{\|x\| \le 1}|(Ax, x)| = \sup_{\|x\| = 1}|(Ax, x)| = \max\left(\sup_{\|x\| = 1}(Ax, x), \, -\!\inf_{\|x\| = 1}(Ax, x)\right).
    \]

    \setcounter{enumi}{0}
    \item Пусть $\lambda \notin [m, M]$. Тогда $d = \mathrm{dist}(\lambda, [m, M]) > 0$. Хотим показать, что $\lambda \in \rho(A)$ --- для этого проверим, что $A_\lambda$ ограничен снизу. Возьмём $x$ с нормой 1: 
    \[
        \|A_\lambda x\| \overset{\text{КБШ}}{\ge} |(A_\lambda x, x)| = |(Ax, x) - \lambda| \ge d.
    \]
    Следовательно, для любого $x$ $\|A_\lambda x\| \ge d \|x\|$, поэтому $\lambda \in \rho(A)$.

    \item Покажем, что $M \in \sigma(A)$. Для этого нам понадобится, что $\|A\| = M$. Если это не так, то применим трюк: рассмотрим оператор $B = A - mI$. Заметим две вещи: во-первых, $\sigma(B) = \sigma(A) - m$, а во-вторых, для $x$ с нормой 1 верно $(Bx, x) = (Ax, x) - m$, поэтому $M_B = M - m$ и $m_B = m - m = 0$. Следовательно, по пункту 3 $\|B\| = M_B$. Если мы докажем, что $M_B \in \sigma(B)$, то получим, что $M = M_B + m \in \sigma(B) + m = \sigma(A)$. То есть достаточно рассмотреть случай $\|A\| = M$.

    Снова проверим критерий принадлежности спектру ССО. Возьмём последовательность $\{x_n\}$, такую что $\|x_n\| = 1$ и $(Ax_n, x_n) \to M$ (найдётся по определению $M$), и посмотрим на $A_M x_n$:
    \[
        \|A_M x_n\|^2 = \|Ax_n\|^2 - (Ax_n, Mx_n) - (M x_n, Ax_n) + M^2\|x_n\|^2 \overset{\text{ССО}}{=} {\underbrace{\|Ax_n\|}_{\le \|A\|}}^2 - 2M\underbrace{(Ax_n, x_n)}_{\to M} + M^2.
    \]
    Так как $\|A\| = M$, получаем $\|A_M x_n\| \to 0$. Значит, $M \in \sigma(A)$.
    
    Осталось проверить, что $m \in \sigma(A)$. Для этого перейдём к $B = -A + MI$. Тогда, опять же, $\sigma(A) = M - \sigma(B)$, $M_B = -m + M$, $m_B = -M + M = 0$. По доказанному $M - m \in \sigma(B)$, и благополучно $m = M - (M - m) \in \sigma(A)$.
\end{enumerate}

\QED



\section{(§12) Компактные операторы}

\subsection{Свойства компактных операторов.}

\textbf{Напоминание.} Подмножество $S$ метрического пространства $E$ называется \textit{предкомпактным}, если его замыкание в $E$ компактно.

\textbf{Замечание.} Полезный критерий: множество $S$ предкомпактно в $E$ тогда и только тогда, когда из любой последовательности в $S$ можно выделить подпоследовательность, сходящуюся в $E$ (но не всегда в $S$). Следует из аналогичного критерия компактности.

\textbf{Определение-утверждение.} Пусть $E_1, E_2$ --- нормированные пространства, $A \in \mathcal L(E_1, E_2)$.
Оператор $A$ называется \textit{компактным}, если выполнено одно из эквивалентных условий:
\begin{enumerate}
    \item образ любого ограниченного множества в $E_1$ является предкомпактным в $E_2$,

    \item образ замкнутого единичного шара $\overline{B}(0, 1)$ предкомпактен,

    \item из любой ограниченной последовательности $\{x_n\}$ можно выделить подпоследовательность $\{x_{n_k}\}$, такую что $\{Ax_{n_k}\}$ сходится в $E_2$.
\end{enumerate}

\textbf{Доказательство.} (эквивалентности определений)

($2 \Rightarrow 1$) Во-первых, подмножество предкомпактного множества предкомпактно (следует из критерия предкомпактности). Во-вторых, любое ограниченное множество $S$ можно поместить в шар $\overline B(0, R)$, $R > 0$. Следовательно, $A(S) \subset A(\overline B(0, R)) = R A(\overline B(0, 1))$, и остаётся проверить предкомпактность последнего. А она есть, поскольку $A(\overline B(0, 1))$ предкомпактно, а умножение на $R$ --- изоморфизм, поэтому сохраняет это свойство.

($1 \Rightarrow 3$) Последовательность ограничена, поэтому её образ $\{y_n\} = \{Ax_n\}$ предкомпактен. По замечанию выше можно выделить подпоследовательность $\{y_{n_k}\} = \{Ax_{n_k}\}$, сходящуюся в $E_2$.

($3 \Rightarrow 2$) Пусть $\{y_n\} = \{Ax_n\}$ --- последовательность, где $x_n \in \overline B(0, 1)$. Собственно, выделим подпоследовательность $\{x_{n_k}\}$, такую что сходится $\{y_{n_k}\} = \{Ax_{n_k}\}$.

\QED

\textbf{Утверждение 1.} Если $\dim(E_1) < \infty$ или $\dim(E_2) < \infty$, то любой $A \in \mathcal L(E_1, E_2)$ компактен.

\textbf{Доказательство.} В обоих случаях образ $A$ конечномерный. Тогда ограниченное множество переходит в конечномерное и ограниченное, то есть предкомпактное.

\QED

\textbf{Утверждение 2.} $I \in \mathcal K(E)$ тогда и только тогда, когда размерность $E$ конечна.

\textbf{Доказательство.} Из \hyperref[th:4-1]{теоремы о компактности единичной сферы} можно получить аналогичный критерий конечномерности пространства: компактность замкнутого шара (так как подмножество вполне ограниченного множества вполне ограничено). Теперь просто используем определение 2: $I(\overline B(0, 1)) = \overline B(0, 1)$, а компактность замкнутого множества совпадает с предкомпактностью.

\QED

\textbf{Утверждение 3.} $\mathcal K(E)$ --- двусторонний идеал в $\mathcal L(E)$. То есть, если $A \in \mathcal K(E)$ и $B \in \mathcal L(E)$, то $AB, BA \in \mathcal K(E)$, а если ещё и $B \in \mathcal K(E)$, то $A + B \in \mathcal K(E)$.

\textbf{Доказательство.} Для $AB$ очевидно: если $S \subset E$ ограничено, то $B(S)$ ограничено, значит, $A(B(S))$ --- предкомпакт.

Для $BA$ воспользуемся определением 3: если выделили подпоследовательность $\{x_{n_k}\}$, такую что $\{Ax_{n_k}\}$ сходится, то $\{BAx_{n_k}\}$ тоже сходится.

Для $A + B$: сначала выделим подпоследовательность, такую что сходится $\{Ax_{n_k}\}$, а потом под-подпоследовательность, чтобы сходилось ещё и $\{Bx_{n_{k_m}}\}$.

\QED

\textbf{Утверждение 4.}(12.4) \phantomsection \label{st:12.4} Если $\dim(E) = \infty$, $A \in \mathcal K(E)$ обратим, то $A^{-1}$ не может быть элементом $\mathcal L(E)$.
Грубо говоря, компактные операторы, действующие из бесконечномерных пространств, ``теряют информацию``.

\textbf{Доказательство.} От противного: $A^{-1} A = I$, $A^{-1} \in \mathcal{L}(E)$.
Тогда по утверждению 3 произведение лежит в $\mathcal K(E)$ --- противоречие с утверждением 2.

\QED

\textbf{Утверждение 5.} (б/д) Если $A \in \mathcal K(E)$, $x_n \xrightarrow{w} x$, то $\|Ax_n - Ax\| \to 0$.
Более того, если $E$ рефлексивно, то верно и обратное (всякий линейный ограниченный оператор, переводящий слабо сходящиеся последовательности в сходящиеся по норме, компактен).

\textbf{Утверждение 6.} (б/д) Если $E$ --- банахово пространство, $A \in \mathcal L(E)$, то $A \in \mathcal K(E)$ тогда и только тогда, когда $A^* \in \mathcal K(E^*)$.

\phantomsection \label{th:12-1} \textbf{Теорема.} (12.1) Пусть $E_1$ --- нормированное пространство, $E_2$ --- банахово пространство, $A_n$ --- последовательность компактных операторов, такая что $\|A_n - A\| \to 0$.
Тогда $A$ компактен.
Другими словами, $\mathcal{K}(E_1, E_2)$ замкнуто в $\mathcal{L}(E_1, E_2)$.

\textbf{Доказательство.} 
На самом деле достаточно того, что для любого $\varepsilon > 0$ найдётся $n$, такое что $\|A_n - A\| < \varepsilon$, то есть $\sup_{\|x\| \le 1} \|A_n x - Ax \| < \varepsilon$.
Раз $E_2$ банахово, то предкомпактность эквивалентна вполне ограниченности.
Далее будем искать $\varepsilon$--сеть во множестве $A (\overline B(0, 1))$.

Из компактности получаем, что $A_n \overline B$ является предкомпактным, то есть вполне ограниченным.
Значит, найдётся $\varepsilon$--сеть $\{y_1^n, \dots, y_m^n\}$ для $A_n \overline B$.
Получается, что для любого $x \in \overline B$ найдётся $y_k^n$, такой что $\|A_n x - y_k^n\| < \varepsilon$.
Итак,
\[
    \min_{1 \le k \le m} \|Ax - y_k^n\| \le \min_{1 \le k \le m} \left( \|Ax - A_n x \| + \|A_n x - y_k^n \| \right) < 2\varepsilon.
\]
То есть $\{y_1^n, \dots, y_m^n\}$ --- 2$\varepsilon$ сеть для $A \overline B$.

\QED

\textbf{Замечание 1.} Поточечной сходимости $A$ недостаточно, и контрпример придумать несложно.

\textbf{Замечание 2.} Данная теорема и определение --- это два самых популярных способа доказательства компактности.
Например, если образы $A_n$ конечномерны и $\|A_n - A\| \to 0$, то $A$ компактен.



\subsection{Свойства собственных значений компактного оператора.}

\phantomsection \label{th:12-2} \textbf{Теорема.} (12.2) Пусть $E(\mathbb C)$ --- банахово пространство, $A \in \mathcal K(E)$, $\lambda \in \mathbb C$, $\lambda \ne 0$.
Тогда $\dim(\Ker A_\lambda) < \infty$.

\textbf{Замечание.} $\lambda = 0$ не работает из-за того, что тогда ядром $A_\lambda$ будет просто ядро $A$, и про него ничего не понятно.

\textbf{Доказательство.} Положим $L := \Ker A_\lambda$. Рассмотрим $\tilde A: L \to L$, $\tilde A := A|_L = \lambda I_L$. Это компактный оператор, так как $A$ компактен, а $L$ замкнуто. Но если $\lambda \neq 0$, то у $\tilde A$ есть ограниченный обратный: $\frac{1}{\lambda}I_L$. Значит, по \hyperref[st:12.4]{утверждению 12.4} $L$ конечномерно.

\QED

\phantomsection \label{th:12-3} \textbf{Теорема.} (12.3) Пусть $E(\mathbb C)$ --- банахово пространство, $A \in \mathcal K(E)$.
Тогда для любого $\delta > 0$ вне круга $\{|\lambda| \le \delta\}$ может находиться лишь конечное число собственных значений оператора $A$.

\textbf{Доказательство.} Будем доказывать в случае гильбертова $E = H$ и самосопряжённого $A$.
От противного: нашлось $\delta_0$, такое что вне круга собственных значений бесконечно много.
Выделим последовательность $\{\lambda_n\}$, тогда каждому из них можно сопоставить единичный собственный вектор $e_n$.
Так как это собственные векторы, для всех $i \ne j$ векторы $e_i$ и $e_j$ ортогональны.

Итак, рассмотрим последовательность $\{Ae_n\}$ --- предкомпакт.
Для любых $i \ne j$:
\[
    \|Ae_i - Ae_j\|^2 = \|Ae_i\|^2 + \|Ae_j\|^2 = |\lambda_i|^2 + |\lambda_j|^2 \ge 2\delta_0^2.
\]
Следовательно, в $\{Ae_n\}$ не удастся выделить сходящуюся подпоследовательность --- противоречие.

\QED

\textbf{Замечание.} В общем случае можно несложно доказать через теорему о почти перпендикуляре, но сейчас просто знакомимся с полезными приёмами.

\textbf{Следствие 1.} Для любого $\delta > 0$ выполнено $\sum_{|\lambda| > \delta} \dim(\Ker A_\lambda)) < \infty$.

\textbf{Следствие 2.} Если $A$ компактен, то точечный спектр $\sigma_p(A)$ не более, чем счётный.



\subsection{Теорема Фредгольма для компактных самосопряжённых операторов.}

\phantomsection \label{th:12-4} \textbf{Теорема.} (Фрéдгольма, 12.4, частный случай)
Пусть $H(\mathbb C)$ --- гильбертово пространство, $A: H \to H$ --- компактный самосопряжённый оператор, $\lambda \in \mathbb C \setminus \{0\}$.
Тогда $\im(A_\lambda) \oplus \Ker(A_\lambda) = H$.

\textbf{Лемма 1.} В условиях теоремы 12.4 если $\lambda \in \sigma(A)$ и $\lambda \ne 0$, то $\lambda \in \sigma_p(A)$.

\textbf{Доказательство.} По \hyperref[th:11-3]{теореме 11.3} $\lambda \in \sigma(A)$ тогда и только тогда, когда найдётся последовательность $\{x_n\}$, такая что $\|x_n\| = 1$ и $A_\lambda x_n \to 0$.
Это значит, что $Ax_n - \lambda x_n \to 0$.
Так как $\{x_n\}$ ограничена и $A$ компактен, можно выделить сходящуюся $A x_{n_k} \to y$.

Вместе с тем, что $Ax_{n_k} - \lambda x_{n_k} \to 0$, получаем, что $\lambda x_{n_k} \to y$.
Значит, $x_{n_k} \to \frac{1}{\lambda} y$ (следовательно, $y \neq 0$).
Возьмём образ $A$: $Ax_{n_k} \to \frac{1}{\lambda} A y$, и в то же время $Ax_{n_k} \to y$ по построению.
Следовательно, $y = \frac{1}{\lambda} Ay$, то есть $y$ --- собственный вектор с собственным значением $\lambda$.

\QED

\textbf{Лемма 2.} \phantomsection \label{l:12.2} (Также известное как утверждение 1 параграфа 11) Если $A$ --- самосопряжённый оператор, $M$ инвариантно относительно $A$, то $M^\bot$ тоже инвариантно относительно $A$.

\textbf{Лемма 3.} В условиях теоремы 12.4 выполнено $\overline{\im A_\lambda} = \im A_\lambda$.

\textbf{Доказательство.} Зафиксируем $\lambda$.
По \hyperref[th:11-2]{теореме 11.2} выполнено $H = \overline{\im A_\lambda} \oplus \Ker A_\lambda$.
Так как ядро $A_\lambda$ инвариантно относительно $A$, его можно взять в качестве $M$ и применить лемму 2, откуда $\overline{\im A_\lambda}$ инвариантно относительно $A$.

Положим $\tilde A = A|_{\overline{\im A_\lambda}}: \overline{\im A_\lambda} \to \overline{\im A_\lambda}$ --- компактный самосопряжённый оператор.
Заметим, что $\lambda$ не является собственным значением $\tilde A$, ведь все возможные собственные векторы лежат в $\Ker A_\lambda$.
Следовательно, по лемме 1 $\lambda \in \rho(\tilde A)$, поэтому $\im \tilde A_\lambda = \overline{\im A_\lambda}$. Но $\im \tilde A_\lambda \subset \im A_\lambda$, значит, $\overline{\im A_\lambda} = \im A_\lambda$.

\QED

\textbf{Доказательство.} (Теоремы Фредгольма) По \hyperref[th:11-2]{теореме 11.2} получаем $H = \overline{\im A_\lambda} \oplus \Ker A_\lambda$.
И по лемме 3 имеем $\overline{\im A_\lambda} = \im A_\lambda$.

\QED

\textbf{Теорема'.} (Альтернатива Фредгольма) В условиях теоремы 12.4 рассмотрим уравнения $A_\lambda x = y$ (1) и $A_\lambda z = 0$ (2).
Тогда либо:
\begin{itemize}
    \item Уравнение (1) имеет решение для всех $y$.
    \item Уравнение (2) имеет нетривиальное решение.
\end{itemize}
Во втором случае уравнение (1) разрешимо тогда и только тогда, когда $y$ ортогонально всем решениям уравнения (2).

\textbf{Доказательство.}
Заметим, что первому случаю соответствует $\im A_\lambda = H$, а второму --- $\Ker A_\lambda \ne 0$.

\QED



\subsection{Теорема Гильберта--Шмидта.}

\phantomsection \label{th:12-5} \textbf{Теорема.} (12.5, Гильберта--Шмидта) Пусть $H(\mathbb C)$ --- сепарабельное гильбертово пространство, $H \neq \{0\}$, $A$ --- компактный самосопряжённый оператор.
Тогда в $H$ найдётся ОНБ, состоящий из собственных векторов оператора $A$.

\textbf{Замечание.} Эта теорема является обобщением теоремы из линейной алгебры о том, что самосопряжённый оператор можно диагонализовать.

\textbf{Лемма.} В условиях теоремы Гильберта--Шмидта у $A$ найдётся собственный вектор.

\textbf{Доказательство.} Если $A = 0$, то точно найдётся. Иначе по \hyperref[th:11-5]{теореме 11.5} имеем $0 \ne \|A\| = \max(|m_-|, |m_+|)$.
Дальше вспоминаем, что $m_-$ и $m_+$ лежат в спектре, поэтому по теореме Фредгольма доказано.

\QED

\textbf{Доказательство.} (теоремы) Для каждого ненулевого собственного значения выберем ортонормированный базис в соответствующем собственном подпространстве и объединим эти базисы. По следствиям из \hyperref[th:12-3]{теоремы 12.3} объединение не более чем счётно: $e_1 \bot e_2 \bot \dots$.
Получаем ортонормированную систему, но это ещё не всё, ибо мы добавили векторы только для ненулевых собственных значений.
Засим рассмотрим базис $\Ker A$ (найдётся, так как пространство сепарабельно) и добавим его к $\{e_n\}$.
Теперь уже докажем, что получился базис, а для этого достаточно доказать, что $\overline {[\{e_n\}]} = H$. Обозначим левую часть за $M$.

По \hyperref[th:4-3]{теореме 4.3} выполнено $H = M \oplus M^\bot$, засим покажем, что $M^\bot = \{0\}$.
Заметим, что $M$ инвариантно относительно $A$ (без замыкания очевидно, а с замыканием тоже, так как $A$ непрерывен).
Следовательно, по \hyperref[l:12.2]{лемме 12.2} $M^\bot$ тоже инвариантно относительно $A$.

Рассмотрим $\tilde A = A|_{M^\bot}\colon M^\bot \to M^\bot$ --- компактный самосопряжённый оператор, так как $M^\bot$ замкнуто.
Допустим, что $M^\bot$ нетривиально. Тогда по предыдущей лемме у $\tilde A$ найдётся собственный вектор, но они все лежат в $M$ --- противоречие.

\QED

\textbf{Следствие.}
Рассмотрим уравнение $A_\lambda x = y$, $\lambda \in \rho(A)$, $A$ --- КССО из теоремы Гильберта-Шмидта.
Тогда можно рассмотреть ОНБ из собственных векторов $\{e_n\}$ и разложить $x = \sum (x, e_n) e_n$ и $y = \sum (y, e_n) e_n$.
Значит, в силу ортогональности $\{e_n\}$ можно записать для всех $n$
\[
    \lambda_n (x, e_n) - \lambda (x, e_n) = (y, e_n) \Rightarrow (x, e_n) = \frac{(y, e_n)}{\lambda_n - \lambda}.
\]



\section{(§13) Элементы нелинейного анализа}

\subsection{Производная Фреше, производная Гато. Формула конечных приращений.}

Пусть $E_1, E_2$ --- банаховы пространства, $D \subset E_1$ открыто, $x_0 \in D$, $F: D \to E_2$.

\textbf{Определение.} $F$ \textit{дифференцируема по Фреше} в точке $x_0$, если найдётся $A \in \mathcal L(E_1, E_2)$ (зависящий от $x_0$), такой что
\[
    F(x_0 + h) - F(x_0) = Ah + o(\|h\|).
\]
$A$ называется \textit{производной (Фреше)} отображения $F$ в точке $x_0$ и обозначается $A = F'(x_0)$.

\textit{Дифференциалом} называется оператор, действующий следующим образом:
\[
    \dif F(x_0, h) = F'(x_0) h = Ah.
\]

\textbf{Утверждение.} 
\begin{enumerate}
    \item Дифференцируемость влечёт непрерывность (в точке).
    \item Если $F \equiv const$, то $F' \equiv 0$.
    \item Если $F \in \mathcal L(E_1, E_2)$, то $F' \equiv F$.
    \item Линейность: $(\alpha F + \beta G)' = \alpha F' + \beta G'$.
\end{enumerate}

\textbf{Определение.} \textit{Производной по направлению} называется 
\[
    \Dif F(x_0, h) = \frac{\dif}{\dif t} F(x_0 + th)\bigg|_{t = 0}.
\]
Если существует $A \in \mathcal L(E_1, E_2)$, такой что $\Dif F(x_0, h) = Ah$ для всех $h$, то $F$ называется \textit{дифференцируемой по Гато} (Gateaux) в точке $x_0$.

\textbf{Утверждение.} Из дифференцируемости по Фреше следует дифференцируемость по Гато и $F_{\text{Фр}}' = F_{\text{Г}}'$, но в обратную сторону это неверно (например, можно взять индикатор графика параболы $\mathbb{R}^2 \to \mathbb{R}$ с выколотой точкой).

\textbf{Теорема.} (13.1, О дифференцируемости композиции) Пусть $X, Y, Z$ --- нормированные пространства, $F: X \to Y$ дифференцируема в $x_0 \in X$, $G: Y \to Z$ --- в точке $y_0 = F(x_0) \in Y$.
Тогда $H := G \circ F$ дифференцируема в точке $x_0$ и $H'(x_0) = G'(F(x_0)) F'(x_0)$.

\textbf{Доказательство.} Заметим, что $F$ и $G$ непрерывны в $x_0$ и $y_0$ соответственно.
Теперь напишем определения дифференцируемости, начиная расписывать, как и полагается в матрёшках, с внешней.
\[
    \Delta z = G'(y_0) \Delta y + \varepsilon_1(\Delta y) \|\Delta y\|,
\]
где $\varepsilon_1 \to 0$ при $\Delta y \to 0$.
\[
    \Delta y = F'(x_0) \Delta x + \varepsilon_2(\Delta x) \|\Delta x\|,
\]
где $\varepsilon_2 \to 0$ при $\Delta x \to 0$.

Теперь мы хотим представить $\Delta z$ в виде $H'(x_0) \Delta x + o(\Delta x)$.
Для этого подставим $\Delta y$ в формулу для $\Delta z$:
\[
    \Delta z = G'(y_0) \big( F'(x_0) \Delta x + \underbrace{\varepsilon_2(\Delta x) \|\Delta x\|}_{(1)} \big) + \underbrace{\varepsilon_1(\Delta y) \|\Delta y\|}_{(2)}.
\]
Заметим, что $(1)$ --- это $o(\Delta x)$, а значит остаётся им после применения $G'(y_0)$ в силу ограниченности оператора.
Теперь $(2)$: разделим на $\|\Delta x\|$ и возьмём норму:
\[
    \frac{\|\varepsilon_1(\Delta y) \| \cdot \|\Delta y\|}{\|\Delta x\|} = \|\varepsilon_1(\Delta y) \| \cdot \frac{\bigg\|F'(x_0) \Delta x + \varepsilon_2(\Delta x) \|\Delta x\| \bigg\|}{\|\Delta x\|} \le
\]
\[
    \le \|\varepsilon_1(\Delta y) \| \cdot \frac{\|F'(x_0) \| \cdot \|\Delta x\| + \|\varepsilon_2(\Delta x) \| \cdot \|\Delta x\|}{\|\Delta x\|}.
\]
Итак, произведение бесконечно малой $\|\varepsilon_1\|$ на ограниченную --- это бесконечно малая, что завершает доказательство.

\QED

\textbf{Теорема.} (13.2, О среднем) Пусть $E_1$, $E_2$ --- вещественные нормированные пространства, $D \subset E_1$ --- выпуклая область, $x_0, x_1 \in D$.
Положим $x(t) = x_0 + t(x_1 - x_0)$ для $t \in [0, 1]$.
Пусть $F: D \to E_2$ дифференцируема.
Тогда
\[
    \|F(x_1) - F(x_0)\| \le \sup_{t \in (0, 1)} \|F'(x(t))\| \cdot \|x_1 - x_0\|.
\]

\textbf{Доказательство.} Положим $\phi(t) = f(F(x(t)))$ для $f \in E_2^*$, который определим позже. Тогда по теореме о композиции $\phi$ дифференцируемо и $\phi'(t) = f(F'(x(t)) \Delta x)$, где $\Delta x = x_1 - x_0$. Применим к $\phi$ теорему Лагранжа:
\[
    |\phi(1) - \phi(0)| \le \|f\| \cdot \sup_{t \in (0, 1)} \|F'(x(t))\| \cdot \|\Delta x\|.
\]
Слева получаем $|f(F(x_1) - F(x_0))|$, поэтому можно по следствию 2 из теоремы Хана--Банаха взять функционал $f$, такой что $\|f\| \le 1$ и $f(F(x_1) - F(x_0)) = \|F(x_1) - F(x_0)\|$.
Заметим, что мы получили ровно то, что нужно.

\QED



\section{(§14) Преобразование Фурье и свертка в пространствах \texorpdfstring{$L_1(\mathbb{R})$ и $L_2(\mathbb{R})$}{L1(R) и L2(R)}}

\subsection{Определения и основные свойства. Формула умножения. Преобразование Фурье свертки.}

В данном параграфе мы не будем различать классы эквивалентности функций в $L_1, L_2$ и их представителей.
На всякий случай: $\hat f'$ --- производная преобразования Фурье, $\widehat{f'}$ --- преобразование Фурье производной.

\textbf{Определение.} Пусть $f \in L_1$.
\textit{Преобразованием Фурье} называется 
\[
    \hat f(y) = (Ff)(y) = \int_{-\infty}^{+\infty} f(x) e^{-ixy} \dif x.
\]

\textbf{Воспоминания.} В гармоническом анализе доказывалось, что:
\begin{itemize}
    \item Если $f$ достаточно хорошая и $f, f' \in L_1$ то $\widehat{f'}(y) = iy \cdot \hat f(y)$.
    \item Если $f, (x \mapsto xf(x)) \in L_1$, то $\hat f'(y) = -i \cdot \widehat{x f(x)}(y)$.
\end{itemize}

\textbf{Утверждение.}
Пусть $f \in L_1$.
Тогда $\hat f \in C(\mathbb R)$ и $\hat f(y) \to 0$ при $y \to \pm \infty$.

\textbf{Доказательство.} Непрерывность напрямую следует из непрерывности по параметру интеграла с параметром, а второе свойство из теоремы Римана об осцилляции.

\QED

Далее мы будем исследовать свойства оператора $F: L_1(\mathbb R) \to C_0(\mathbb R)$, где $C_0(\mathbb R)$ --- это непрерывные функции, стремящиеся к нулю на бесконечности, с нормой $\sup_{x \in \mathbb R} |f(x)|$.

\textbf{Утверждение.}
$F$ --- линейный ограниченный оператор. 

\textbf{Доказательство.} Линейность очевидна. Ограниченность тоже:
\[
    \|\hat f\|_{C_0(\mathbb R)} = \sup_{y \in \mathbb R} \left|\int_{-\infty}^{+\infty} f(x) e^{-ixy} \dif x\right| \le \int_{-\infty}^{+\infty} |f(x)| \dif x = \|f\|_{L_1(\mathbb R)}.
\]

\textbf{Замечание.} Взяв функцию $\mathbf I_{[-\pi, \pi]}$, можно проверить, что $\|F\| = 1$.

\QED

\textbf{Утверждение.} (формула умножения) Пусть $f, g \in L_1$, тогда
\[
    \int_{\mathbb R} \hat f(x) g(x) \dif x = \int_{\mathbb R} f(x) \hat g(x) \dif x.
\]

\textbf{Доказательство.} Можно применить теорему Фубини, но для этого нужно проверить, что хотя бы один из повторных интегралов конечен. Сделаем это:
\[
    \int_{\mathbb R} \int_{\mathbb R} |f(y)e^{-iyx} g(x)| \dif y \dif x \le \int_{\mathbb R} \int_{\mathbb R} |f(y)||g(x)| \dif y \dif x = \|f\|_{L_1} \|g\|_{L_1}.
\]
И, собственно, по теореме Фубини:
\[
    \int_{\mathbb R} \hat f(x) g(x) \dif x = \int_{\mathbb R} \int_{\mathbb R} f(y)e^{-iyx} g(x) \dif y \dif x =  \int_{\mathbb R} \int_{\mathbb R} f(y)e^{-iyx} g(x) \dif x \dif y = \int_{\mathbb R} f(y) \hat g(y) \dif y.
\]

\QED

\textbf{Определение.} Пусть $f, g \in L_1(\mathbb R)$. \textit{Свёрткой} называется
\[
    (f * g)(x) \overset{\text{п.в.}}{=} \int_{-\infty}^{+\infty} f(\xi) g(x - \xi) \dif \xi.
\]

\textbf{Утверждение.} $f * g \in L_1(\mathbb R)$.
Очевидно, теорема Фубини.

\textbf{Утверждение.} $\widehat{f * g} = \hat f \cdot \hat g$.
Очевидно, теорема Фубини.

\textbf{Утверждение.} Для почти всех $x$ выполнено $(f * g)(x) = (g * f)(x)$.
Очевидно, замена переменной в интеграле.

Далее перейдём к преобразованию Фурье в $L_2(\mathbb R)$, а $F$ домножим для удобства на коэффициент $\frac{1}{\sqrt{2\pi}}$. Проблема в том, что формулу для $L_1$ не получится распространить на $L_2$.
Засим тут имеются разные способы определения:
\begin{enumerate}
    \item (для любознательных) Как предел при $R \to +\infty$ интеграла
        \[
            \frac{1}{\sqrt{2\pi}} \int_{-R}^R f(x) e^{-ixy} \dif x,
        \]
        который сходится к некоторой функции из $L_2$ по теореме Планшереля. Однако этот способ не очень идиоматичен с точки зрения функционального анализа.

    \item Вспомним, что существует пространство Шварца $S$ функций из $C^\infty(\mathbb R)$, у которых все степени производных убывают на бесконечности быстрее любого многочлена. Заметим, что $S$ является линейным многообразием в $L_2$, а также $\overline S = L_2$ (б/д). Также известно, что на $S$ оператор $F: S \to S$ является изоморфизмом (с нормой $L_2$). Тогда по \hyperref[th:5-3]{теореме 5.3} его можно продолжить на $L_2$.
\end{enumerate}

Мы пойдём по второму пути, засим явно напишем все необходимые свойства преобразования Фурье на $S$ из гармонического анализа:

\textbf{Лемма 1.} (б/д) $F: S \to S$ --- это биекция, причём обратный оператор имеет вид
\[
    (F^{-1} g)(x) = \frac{1}{\sqrt{2\pi}} \int_{-\infty}^{+\infty} g(y) e^{ixy} \dif y.
\]

\textbf{Лемма 2.} Для всех $f, g \in S$ выполнено $(f, g) = (\hat f, \hat g)$. Следовательно, $F \in \mathcal{L}(S)$.

\textbf{Доказательство.} Воспользуемся формулой умножения и формулой обращения:
\[
    (\hat f, \hat g) = \int_{\mathbb R} \hat f(x) \overline{\hat g(x)} \dif x = \int_{\mathbb R} f(x) \widehat{\overline{\hat g(x)}} \dif x.
\]
Заметим, что
\[
    \widehat{\overline{\hat g(x)}} = \frac{1}{\sqrt{2\pi}} \int_{\mathbb R} \overline{\hat g(y)} e^{-ixy} \dif y = \overline{\frac{1}{\sqrt{2\pi}} \int_{\mathbb R} \hat g(y) e^{ixy} \dif y} = \overline{g(x)}.
\]
Следовательно,
\[
    (\hat f, \hat g) = \int_{\mathbb R} f(x) \overline{g(x)} \dif x = (f, g).
\]

\QED

\textbf{Теорема.} Преобразование Фурье $F$ в $L_2$ --- это изометрический изоморфизм гильбертовых пространств (сохраняет скалярное произведение).

\textbf{Доказательство.} По лемме 2 $F$ можно продлить до $\overline S = L_2$. По лемме 1 $F^{-1}$ тоже можно продлить до $L_2$. Заметим, что продолжения обратны друг другу: если $\{f_n\} \subset S$ и $f_n \overset{L_2}{\to} f$, то
\[
    F^{-1}Ff = \lim_{n \to \infty} F^{-1}Ff_n = \lim_{n \to \infty} f_n = f.
\]
Аналогично для $FF^{-1}g$.

Со скалярным произведением довольно очевидно: берём $f_n \to f$ и $g_n \to g$ из $S$, тогда $(\hat f_n, \hat g_n) = (f_n, g_n)$, и далее по непрерывности скалярного произведения.

\QED

\textbf{Следствие.} По \hyperref[st:12.4]{утверждению 12.4} $F$ не компактный оператор.

\textbf{Свойство 1.} $(F \phi)(y) = (F^{-1}\phi)(-y)$. Следовательно, $(F^2 \phi)(y) = \phi(-y)$. Следовательно, $F^4 = I$.

\textbf{Доказательство.} Это фактически равенство двух операторов. На $S$ оно выполняется по лемме 1, а значит, продолжается и до замыкания, то есть $L_2$.

\QED

\textbf{Свойство 2.} $F^* = F^{-1}$.

\textbf{Доказательство.} Обозначим $F^{-1}g := \tilde g$. Заметим, что $\hat g(x) = \tilde g(-x)$. Тогда из сохранения скалярного произведения:
\[
    (\hat f, g) = (\hat{\hat{f}}, \hat g) = \int_{\mathbb R} \tilde{\hat{f}}(-x) \overline{\hat g(x)} \dif x = \int_{\mathbb R} f(-x) \overline{\tilde g(-x)} \dif x = \int_{\mathbb R} f(x) \overline{\tilde g(x)} \dif x = (f, \tilde g).
\]

\QED

\textbf{Следствие.} $F$ не самосопряжённый, так как $F^{-1} \neq F$ (например, $\widehat{\mathrm I_{[0, 1]}} \neq \widetilde{\mathrm I_{[0, 1]}}$).

\textbf{Доказательство.} Обозначим $F^{-1}g := \tilde g$. Заметим, что $\hat g(x) = \tilde g(-x)$. Тогда из сохранения скалярного произведения:
\[
    (\hat f, g) = (\hat{\hat{f}}, \hat g) = \int_{\mathbb R} \tilde{\hat{f}}(-x) \overline{\hat g(x)} \dif x = \int_{\mathbb R} f(-x) \overline{\tilde g(-x)} \dif x = \int_{\mathbb R} f(x) \overline{\tilde g(x)} \dif x = (f, \tilde g).
\]

\QED

\textbf{Свойство 3.} $\sigma(F) = \{\pm 1, \pm i\}$.

\textbf{Доказательство.} 
Заметим, что $\sigma(F) \subset \{\pm 1, \pm i\}$, поскольку $\sigma^4(F) \subset \sigma(F^4) = \sigma(I)$ (это было в теореме про спектральный радиус).

В обратную сторону следует из того, что $\{H_n(x) e^{-\frac{x^2}{2}}\}$ --- это собственные векторы с собственными значениями $(-i)^n$, где $H_n$ --- многочлены Эрмита.
Подробнее в Колмогорове--Фомине.

\QED



\subsection{Операторы Гильберта --- Шмидта.}

\textbf{Определение.} \textit{Оператором Вольтера} называется $V: L_2[0, 1] \to L_2[0, 1]$, 
\[
    (Vf)(x) = \int_0^x f(t) \dif t.
\]
На подынтегральное можно смотреть, как на произведение $f(t)$ и $K(x, t)$, где $K$ --- ядро оператора, индикатор $t \le x$.

\textbf{Определение.} Оператор $A: L_2[a, b] \to L_2[a, b]$ называется \textit{оператором Гильберта--Шмидта}, если существует $K \in L_2([a, b]^2)$, такой что
\[
    (Af)(x) = \int_a^b K(x, t) f(t) \dif t.
\]

\textbf{Замечание.} Стоит доказать корректность определения, то есть что образ лежит в $L_2$, но это будет заметно из доказательства ниже.

\textbf{Теорема.} Все операторы Гильберта--Шмидта компактны.

\textbf{Доказательство.} Как известно, в $L_2([a, b])$ есть ОНБ $\{\phi_i\}$. Тогда в $L_2([a, b]^2)$ есть ОНБ $\{\phi_i(x) \cdot \phi_j(t)\}$ (подробнее в Колмогорове--Фомине).
Разложим $K$ по нему: 
\[
    K(x, t) \overset{L_2}{=} \sum_{i,j} a_{i,j} \phi_i(x) \phi_j(t).
\]
Возьмём конечные суммы:
\[
    K_n(x, t) = \sum_{i,j=1}^n a_{i,j} \phi_i(x) \phi_j(t)
\]
Пусть теперь $A_n$ --- операторы Гильберта-Шмидта с ядром $K_n$. Заметим, что $A_n$ компактные, так как имеют конечномерный образ:
\[
    \int_a^b K_n(x, t) f(t) dt = \sum_{i=1}^n \phi_i(x) \int_a^b \sum_{j=1}^n a_{ij} \phi_j(t) f(t) dt = \sum_{i=1}^n c_i \phi_i(x).
\]
Проверим, что $\|A_n - A\| \to 0$. Обозначим $R_n = K_n - K$.  Нужно просто применить неравенство Коши–Буняковского:
\begin{multline*}
    \|(A_n - A)f\|_{L_2}^2 = \int_{a}^{b} \left|\int_a^b R_n(x, t) f(t) dt \right|^2 dx \le \\ \le \int_{a}^{b} \left(\int_a^b |R_n(x, t)|^2 dt\right) \left(\int_a^b |f(t)|^2 dt\right) dx = \|R_n\|_{L_2}^2 \|f\|_{L_2}^2.
\end{multline*}
$\|R_n\|_{L_2} \to 0$ по определению базиса, значит, $\|A_n - A\|_{\mathcal L(L_2)} \to 0$. Следовательно, $A$ компактен, как предел компактных.

\QED

\end{document}